% Part: proof-theory
% Chapter: normalization
% Section: introduction

\documentclass[../../../include/open-logic-section]{subfiles}

\begin{document}

\olfileid{pt}{nor}{seg}

\olsection{برخې او پرېکونونه}

د !!^a{introduction} قاعدې ورپسې !!a{elimination} قاعده په
!!a{proof} کښې يو تاو راتاو دے، او د !!a{proof} عادي کول د
داسې تاو راتاوونو لرې کول دي۔ خو \Elim{\lor} او
\Elim{\lexists} د ``هغه تاو راتاو چې لرې کول ئې غواړو''
تعريف هم لږ پېچلے کوي۔ ستونزه دا ده چې د دغو قواعدو له امله
په تاو راتاو کښې !!{elimination} قاعده ضرور نۀ ده چې له
!!{introduction} قاعدې \emph{سمدستي} وروسته راشي۔ په منځ
کښې يو \Elim{\lor} يا \Elim{\lexists} راتلے شي، مثلاً:
\[
\AxiomC{}
\DeduceC{$\lexists[x][!C(x)]$}
\AxiomC{$\Discharge{!A}{x}$}
\DeduceC{$!B$}
\DischargeRule{\Intro{\lif}}{x}
\UnaryInfC{$!A \lif !B$}
\DischargeRule{\Elim{\lexists}}{y}
\BinaryInfC{$!A \lif !B$}
\AxiomC{}
\DeduceC{$!A$}
\RightLabel{\Elim{\lif}}
\BinaryInfC{$!B$}
\DisplayProof
\]
په حقيقت کښې هر څومره \Elim{\lor} يا \Elim{\lexists}
کولے شي \Intro{\lif} له \Elim{\lif} څخه بېل کړي، مثلاً:
\[
\AxiomC{}
\DeduceC{$\lexists[x][!C(x)]$}
\AxiomC{}
\DeduceC{$!D \lor !E$}
\AxiomC{$\Discharge{!A}{x}$}
\DeduceC{$!B$}
\DischargeRule{\Intro{\lif}}{x}
\UnaryInfC{$!A \lif !B$}
\AxiomC{$\Discharge{!A}{x}$}
\DeduceC{$!B$}
\DischargeRule{\Intro{\lif}}{x}
\UnaryInfC{$!A \lif !B$}
\RightLabel{\Elim{\lor}}
\TrinaryInfC{$!A \lif !B$}
\DischargeRule{\Elim{\lexists}}{y}
\BinaryInfC{$!A \lif !B$}
\AxiomC{}
\DeduceC{$!A$}
\RightLabel{\Elim{\lif}}
\BinaryInfC{$!B$}
\DisplayProof
\]
دا بېلګه ښيي چې يو~\Elim{\lif} په څو تاو راتاوونو کښې
راتلے شي، ځکه له يوه څخه د زياتو~\Intro{\lif} وروسته دے۔
د دغو ټولو امکانونو د رانغاړلو دپاره تاو راتاوونه د
!!{formula} د موردونو په هغو لړيو تعريفوو چې په
\Log{N1i}-ثبوت کښې د \emph{برخو} په نوم يادېږي۔ زموږ په
بېلګه کښې د~$!A \lif !B$ د موردونو يوه داسې لړۍ ده چې د
$\Elim{\lor}$ يا $\Elim{\lexists}$ له کوچنۍ مقدمې وروسته
د هماغه استنتاج پايله راځي۔ په دې بېلګه کښې دوه داسې برخې
شته چې هره يوه د~$!A \lif !B$ درې موردونه لري۔ اوس ئې کره
تعريف وړاندې کوو۔
\textbf{د سرچينه‌يي ترتيب سمون OLCMP-175:} په تشريحي جمله
کښې د کوچنۍ مقدمې او پايلې ترتيب سرچپه و؛ د لاندې تعريف له
درېيم شرط سره برابر شو۔ رسمي رسمونه او تعريف عين دي۔

\begin{defn}
په \Log{N1i} !!{proof}~$\delta$ کښې \emph{برخه} د هماغه
!!{formula}~$!A$ د موردونو يوه لړۍ $!A_1$، \dots، $!A_n$
ده چې په $\delta$ کښې لاندې شرطونه پوره کوي:
\begin{enumerate}
  \item $!A_1$ د \Elim{\lor} يا \Elim{\lexists} پايله نۀ ده؛
  \item $!A_n$ د \Elim{\lor} يا \Elim{\lexists} کوچنۍ مقدمه نۀ ده؛
  \item د $i = 1$، \dots،~$n-1$ دپاره، $!A_i$ د
  \Elim{\lor} يا \Elim{\lexists} کوچنۍ مقدمه ده او $A_{i+1}$
  د هماغه استنتاج پايله ده۔
\end{enumerate}
د دغې برخې \emph{اوږدوالے}~$n$ دے۔
\end{defn}

هره برخه د لرې کولو هدف، يعنې تاو راتاو، نۀ ده۔ هغه
برخې چې داسې هدفونه وي، د \emph{پرېکونونو} په نوم يادېږي۔

\begin{defn}
\emph{پرېکون} هغه برخه ده چې $!A_n$ د !!a{elimination}
استنتاج لويه مقدمه وي، او يا $n=1$ او $!A_1 = !A_n$ د
!!a{introduction} استنتاج يا~\FalseInt پايله وي، يا~$n>1$
وي۔
\end{defn}

پام وکړئ چې د پرېکون برخه ضرور نۀ ده چې د
!!a{introduction} قاعدې په پايله پيل شي۔ لازمه دا ده چې
د !!a{elimination} قاعدې په لويه مقدمه پاے ته ورسېږي۔ خو
د~$1$ اوږدوالي برخه هغه وخت پرېکون ګڼل کېږي چې د
!!{introduction} قاعدې يا د دروغو ايستلو پايله وي۔
\textbf{د سرچينه‌يي شرط سمون OLCMP-176:} په دې يادونه کښې
د دروغو ايستلو استثنا پاتې وه؛ د پاسني اصلي تعريف همدغه
استثنا په نثر کښې هم وساتل شوه۔

\begin{defn}
د پرېکون \emph{رتبه}~$\depth{!A}$ ده۔ د
!!a{proof}~$\delta$ د \emph{پرېکون رتبه}~$\cutrank{\delta}$
په~$\delta$ کښې د پرېکونونو تر ټولو لوړه رتبه ده۔ پرېکون په
~$\delta$ کښې د \emph{تر ټولو لوړې رتبې} بلل کېږي که رتبه
ئې~$\cutrank{\delta}$ وي، يعنې تر ټولو لوړه رتبه ولري۔
د~$\delta$ \emph{پرېکون اوږدوالے} د هغې د تر ټولو لوړې رتبې
د ټولو پرېکونونو د اوږدوالو مجموعه ده۔
\end{defn}

\end{document}
